\section{Interface Table: TFPT Paper 2 \texorpdfstring{$\leftrightarrow$}{<->} Lean Status}
\label{sec:interface-table}
% ---------------------------------------------------------------- %

The following table maps each load-bearing claim of TFPT
Paper~2 to its current Lean status and the still-missing
upstream proof. It is the primary defence against scope-creep:
the present note formalises columns~1 and~2; column~3 is
explicitly out of scope.

\smallskip
\begin{center}
\footnotesize
\begin{tabularx}{\textwidth}{@{}p{0.30\textwidth}p{0.30\textwidth}Y@{}}
\toprule
\textbf{TFPT Paper 2 claim} & \textbf{Lean status (v11)} & \textbf{Missing upstream proof} \\
\midrule
$\varepsilon^2 = \mathbf 1$ on $E$ &
  Two structural layers: \texttt{CalderonProjector}
  (idempotent $\pi$ gives $\varepsilon := 2\pi-1$ with $\varepsilon^2 = 1$)
  and \texttt{BoundaryPolarization} (any
  $H = H_+ \oplus H_-$ gives such an idempotent via
  \texttt{IsCompl.projection}); feeds
  \texttt{CalderonCertificate}. &
  Analytic construction of a boundary spectral decomposition
  $H = H_+ \oplus H_-$ from Paper 1's elliptic Calder\'on
  operator on Sobolev spaces. \\[3pt]
$\dim E_+ = 2$ &
  Stages A, B, and B+ all formalised
  (\texttt{HiggsIndexCertificate},
  \texttt{HiggsTopForm},
  \texttt{HiggsSchemeCohomologyShadow});
  \texttt{DegreeOneTwoVarForm}\,$K$ $\simeq_{\ell K}$
  \texttt{MvPolynomial.homogeneousSubmodule (Fin 2) K 1}
  $= K[X_0, X_1]_1$, the algebraic-geometry "Proj-side"
  presentation of $H^0(\mathbb P^1, \mathcal O(1))$ &
  Scheme-theoretic identification
  $H^0(\mathrm{Proj}\,K[X,Y], \mathcal O(1)) \simeq K[X, Y]_1$
  (a general algebraic-geometry theorem, independent of TFPT). \\[3pt]
$\dim E_- = 3$ &
  Stage A--D forward + Stage D backward
  (\texttt{YukawaStageDExistence}, v10):
  $\dim E_- = 3 \Rightarrow \exists\,\omega$
  with \texttt{ContractionInjective} and
  \texttt{ContractionSurjective}, via $b.\texttt{det}$ and
  direct $3{\times}3$ matrix computation. The Stage D
  characterisation
  $\exists\,\omega\;(\textrm{CI} \wedge \textrm{CS})
    \iff \dim E_- = 3$ holds under \texttt{Nontrivial}\,$E_-$. &
  TFPT-internal construction of $\omega$ from boundary kernel
  and representation data (the existence is now algebraically
  closed by basis-determinant; what remains is whether TFPT's
  boundary derivation produces \emph{this particular}
  $\omega$). \\[3pt]
$3 q_- + 2 q_+ = 0$ &
  Theorem from determinant character of carrier torus
  (\texttt{DeterminantCharacter}) &
  Seam-winding $[u_\Sigma] = 1$ produces the determinant-
  preserving condition. \\[3pt]
$(q_-, q_+) = (-2, 3)$ &
  Theorem from general $(m,n)$ lattice rigidity
  (\texttt{LatticeRigidityGeneral}) with
  \texttt{PrimitiveOrientedDeterminantCarrier} as input &
  Orientation $q_- < 0 < q_+$ and primitivity
  $\gcd = 1$ from seam-class data. \\[3pt]
$Y = \mathrm{diag}(-\tfrac13,-\tfrac13,-\tfrac13,\tfrac12,\tfrac12)$ &
  Bundled theorem
  (\texttt{CarrierPremises.hypercharge\_carrier\_packet}) plus
  concrete $5\times 5$ model (\texttt{Hypercharge}) &
  None (this is the conclusion). \\[3pt]
$\tr_E Y = 0$ &
  Bundled theorem
  (\texttt{CarrierPremises.trace\_Y\_eq\_zero}) &
  None. \\[3pt]
$6\, Y^2 - Y - \mathbf 1 = \mathbf 0$ &
  Bundled theorem
  (\texttt{CarrierPremises.carrier\_polynomial\_Y}) in
  $\mathrm{End}_K M$ &
  None. \\[3pt]
$S^+ = \Lambda^{\mathrm{even}} E$, family $\nu^c$ &
  \emph{Not formalised} &
  Exterior-power packet construction; chiral-family decomposition. \\[3pt]
$N_{\mathrm{fam}} = 3$ &
  \emph{Not formalised} &
  Seam winding $\to$ family count. \\[3pt]
$G_{\mathrm{phys}} = (SU(3)\times SU(2)\times U(1)_Y)/\mathbb Z_6$ &
  \emph{Not formalised} &
  Stabilizer theorem for the carrier datum; $\mathbb Z_6$
  quotient theorem. \\[3pt]
SM hypercharges on all chiral multiplets &
  \emph{Not formalised} &
  Derived representations: dual, exterior powers,
  tensor products of the carrier. \\[3pt]
$b_1 = 41/10$, $\Omega_{\mathrm{adm}} = 48$, $N_\Phi = 1$ &
  \emph{Not formalised} &
  Index ledger; representation audit. \\
\bottomrule
\end{tabularx}
\end{center}

\smallskip\noindent
The pattern is clear. The Lean note formalises the
\emph{carrier-rigidity algebra} of Paper 2 in full; the
\emph{TFPT-internal derivations} feeding that algebra remain
upstream work; and the \emph{downstream SM-packet content}
(derived representations, $\mathbb Z_6$ quotient, family
counting, index ledger) is intentionally out of scope.

% ---------------------------------------------------------------- %
\section{Limitations}
\label{sec:limits}
% ---------------------------------------------------------------- %

\paragraph{Yukawa: Stage D resolves the v5 critique.}
The v5 critique of Stage C was that it bundled the contraction
iso as data, so it was equivalent (in finite dimensions) to
$\dim E = 3$ and did not "explain" the rank. v6 resolved this
by introducing Stage D; v7 renamed its operational predicates
to keep semantic-load words out of the Lean interface:

\[
\underbrace{\omega : \Lambda^3 E_- \to K}_{\text{primary datum}}
\;\;
\underbrace{+\;\texttt{ContractionInjective}\;+\;\texttt{ContractionSurjective}\;+\;\texttt{Nontrivial}\,E_-}_{\text{operational predicates and non-triviality}}
\;\xRightarrow{\text{Lean theorem}}\;
\dim E_- = 3.
\]

The contraction $C_\omega$ is a \emph{constructed} linear map;
the iso conclusion is a Lean theorem under the two predicates
and non-triviality of $E_-$; the rank-3 statement follows from
the iso via elementary dimension counting. There is no longer
any structural ambiguity at the Yukawa interface in Lean.

\paragraph{Why \texttt{Nontrivial}\,$E_-$ is necessary.}
Without the non-triviality hypothesis the rank-3 slogan would be
false: for $E_- = 0$ the spaces $\Lambda^2 E_-$ and $\Lambda^3
E_-$ are both zero, the contraction $C_\omega : 0 \to 0$ is
trivially bijective, both operational predicates hold vacuously,
but $\dim E_- = 0 \ne 3$. The transparency theorem
\texttt{finrank\_E\_eq\_zero\_or\_three} (v9) records the
unconditional disjunction $\dim E_- = 0 \lor \dim E_- = 3$;
the rank-3 conclusion is the case under
\texttt{Nontrivial}\,$E_-$. Both signatures are signature-locked
in \texttt{AuditContract.lean}.

\paragraph{Stage D characterisation, both directions formalised (v10).}
For all four Yukawa stages the note now proves the full
characterisation
$\text{Nonempty}\,(\text{Stage}) \iff \dim E = 3$.
For Stage D, under \texttt{Nontrivial}\,$E$:
\[
\text{Nonempty}\,\bigl\{\,Y : \texttt{YukawaTrilinearForm}\,K\,E
  \;\big|\; Y.\texttt{ContractionInjective}
  \land Y.\texttt{ContractionSurjective}\,\bigr\}
\;\iff\;
\dim E = 3.
\]
The forward direction is
\texttt{finrank\_E\_eq\_three\_of\_condition}; the backward
direction is \texttt{YukawaStageDExistence.exists\_yukawaTrilinearForm\_of\_finrank\_eq\_three},
proved via \texttt{Module.Basis.det} transported through
\texttt{exteriorPower.alternatingMapLinearEquiv} and direct
$3\times 3$ matrix-determinant computation. The combined
characterisation is \texttt{stage\_D\_iff\_finrank\_eq\_three}.

\paragraph{What remains is upstream TFPT theory work.}
The Lean chain takes the Stage D quadruple
$(\omega, \texttt{ContractionInjective},
\texttt{ContractionSurjective}, \texttt{Nontrivial}\,E_-)$
as input. The remaining open question is whether TFPT itself
produces such a quadruple from the boundary kernel. Concretely:

\begin{itemize}[leftmargin=1.6em, itemsep=1pt]
\item \emph{Existence of $\omega$.}\ Paper 2's Yukawa-type
    argument should produce an explicit alternating trilinear
    form on $E_-$. This is a representation-theoretic statement
    about the boundary kernel output, not a Lean engineering
    statement.
\item \emph{Verification of the two predicates.}\ Once $\omega$
    is constructed, the two predicates need to be checked.
    These checks are local algebraic statements (injectivity:
    no nonzero $v$ vanishes against all wedges; surjectivity:
    every dual functional on $\Lambda^2 E_-$ is realised as a
    contraction).
\end{itemize}

Neither task requires modifying the present Lean chain. They are
TFPT mathematical work, properly upstream of this experiment.

\paragraph{Remaining limitations of the Lean chain itself.}
\begin{itemize}[leftmargin=1.6em, itemsep=1pt]
\item \textbf{Calder\'on: two structural layers formalised,
    analytic input of the polarisation open.}\
    \texttt{CalderonProjector} proves $\pi^2 = \pi
    \Rightarrow \varepsilon^2 = \mathbf 1$, and
    \texttt{BoundaryPolarization} (v10) shows that any
    direct-sum decomposition $H = H_+ \oplus H_-$ of a
    $K$-module produces such an idempotent via Mathlib's
    \texttt{IsCompl.projection}. What remains is the analytic
    Paper 1 theorem producing the boundary spectral
    decomposition itself --- a question of boundary elliptic
    theory on Sobolev spaces, independent of the carrier
    algebra.
\item \textbf{Higgs: algebraic Stage B and scheme-cohomology
    shadow formalised, Proj construction open.}\
    \texttt{HiggsTopForm} identifies $E_+$ with the algebraic
    shadow of $H^0(\mathbb P^1, \mathcal O(1))$ (degree-1 binary
    forms), and v10's \texttt{HiggsSchemeCohomologyShadow}
    refines this further: the algebraic shadow is now exhibited
    as a Mathlib-native object,
    \texttt{MvPolynomial.homogeneousSubmodule (Fin 2) K 1}.
    What remains is the general algebraic-geometry theorem
    $H^0(\mathrm{Proj}\,K[X,Y], \mathcal O(1)) \simeq K[X, Y]_1$,
    which lives in \texttt{Mathlib.AlgebraicGeometry.*} and is
    independent of TFPT.
\item \textbf{Yukawa: Stage D both directions formalised.}\
    v10's \texttt{YukawaStageDExistence} closes the backward
    direction:
    $\dim E_- = 3 \Rightarrow \exists\,\omega$ with
    \texttt{ContractionInjective} and
    \texttt{ContractionSurjective}, via $b.\texttt{det}$ and
    direct $3{\times}3$ matrix-determinant computation. The
    characterisation
    $\exists\,\omega\;(\textrm{CI} \wedge \textrm{CS})
      \iff \dim E_- = 3$ holds under
    $\texttt{Nontrivial}\,E_-$. What remains is the
    TFPT-internal question: does the boundary kernel produce
    \emph{the same} $\omega$ (or one equivalent up to action of
    the carrier symmetries)?
\item \textbf{Scope discipline.}\ The Lean chain formalises the
    $3+2$ carrier hypercharge. It does \emph{not} formalise the
    full SM hypercharge spectrum on duals, exterior powers, and
    tensor products. The $\mathbb Z_6$ quotient and the
    derived-representation extensions are out of scope. See
    \S\ref{sec:next}.
\end{itemize}

\paragraph{Other items deliberately out of scope.}
\begin{itemize}[leftmargin=1.6em, itemsep=0pt]
\item Full $\mathbb Z_6$-quotient theorem for
    $G_{\mathrm{phys}} = (SU(3)\times SU(2)\times U(1)_Y)/\mathbb Z_6$.
\item Derived Standard-Model representations (dual, exterior
    powers, tensor products); the diagonal formalised here is
    the carrier hypercharge, not the full SM hypercharge
    spectrum.
\item Analytic content of Paper~1 (Calder\'on polarisation,
    boundary elliptic theory).
\end{itemize}

% ---------------------------------------------------------------- %
\section{How It Was Verified}
% ---------------------------------------------------------------- %

The audit gate \texttt{scripts/audit.sh} enforces seven hard
checks on every commit:

\begin{enumerate}[leftmargin=1.6em, itemsep=0pt]
\item \texttt{lake build} succeeds (currently $\sim 1\,800$ jobs,
      $\sim 10$~s on top of the Mathlib cache).
\item No \texttt{sorry} or \texttt{admit} in the source.
\item No \texttt{axiom} or \texttt{constant} declarations.
\item No \texttt{unsafe}, \texttt{opaque}, or \texttt{partial}
      declarations.
\item \texttt{\#print axioms} on every headline theorem returns a
      subset of $\{\texttt{propext},\ \texttt{Classical.choice},\
      \texttt{Quot.sound}\}$.
\item Every headline name in
      \texttt{TfptCarrier/AuditCheck.lean} elaborates via
      Lean's \texttt{\#check}. This catches renames, deletions,
      and elaboration failures that a text search would miss.
\item Every headline theorem in
      \texttt{TfptCarrier/AuditContract.lean} elaborates inside
      an
      \texttt{example : <exact intended type> := <theorem>}
      pattern. This is strictly stronger than \texttt{\#check}:
      if a future refactor silently weakens a theorem's
      conclusion without renaming it, the
      \texttt{example} fails to typecheck and the audit fails.
\end{enumerate}

Current verbatim output of \texttt{audit.sh}:

\begin{lstlisting}
=== Build ===
  [PASS] lake build succeeded

=== sorry / admit ===
  [PASS] no sorry or admit found

=== axiom / constant declarations ===
  [PASS] no axiom or constant declarations

=== unsafe / opaque / partial declarations ===
  [PASS] no unsafe / opaque / partial declarations

=== axiom audit (#print axioms) ===
  [PASS] all theorems use only [propext, Classical.choice, Quot.sound]

=== headline theorems (Lean #check) ===
  [PASS] AuditCheck.lean elaborated; all headline #checks succeeded

=== theorem signatures (Lean example signature-lock) ===
  [PASS] AuditContract.lean elaborated; all signature locks hold

=== Final verdict ===
  AUDIT: PASS
\end{lstlisting}

\texttt{\#eval} smoke tests in \texttt{Sanity.lean}:

\begin{lstlisting}
info: TfptCarrier/Sanity.lean:27:0: 0
info: TfptCarrier/Sanity.lean:30:0:
  !![0, 0, 0, 0, 0; 0, 0, 0, 0, 0; 0, 0, 0, 0, 0;
     0, 0, 0, 0, 0; 0, 0, 0, 0, 0]
\end{lstlisting}

% ---------------------------------------------------------------- %
\section{Changelog}
\label{sec:changes}
% ---------------------------------------------------------------- %

